% One in seventy: a determinism audit of an OCR ensemble
%
% Deliberately plain LaTeX. Only amsmath, booktabs, hyperref, url and graphicx,
% all of which are in every TeX distribution and on arXiv's. Nothing here needs a
% custom class, a bibliography style file or a package that might be missing on a
% submission server, because the cost of a build failure at submission time is far
% higher than the cost of a plainer document.
\documentclass[11pt,a4paper]{article}

\usepackage[utf8]{inputenc}
\usepackage[T1]{fontenc}
\usepackage{amsmath}
\usepackage{booktabs}
\usepackage{url}
\usepackage[hidelinks]{hyperref}
\usepackage[margin=1in]{geometry}

\title{One in seventy decisions depended on dictionary insertion order:\\
a determinism audit of a multi-engine OCR ensemble}

\author{Igor Saevets\\
\small\url{https://github.com/igorsaevets/scanquorum}}

\date{August 2026}

\begin{document}
\maketitle

\begin{abstract}
Combining several OCR engines by voting is more than thirty years old. We report a
permutation-order determinism audit of one such system. We have found no prior report of
this test on an OCR voting ensemble, and we do not claim priority.
We generated a fixture of 11{,}780 distinct decisions from a running ensemble, replayed
each one under every permutation of the order in which the engines are supplied, and
required the set of outcomes to have size one. It did not. On 165 decisions, or 1.40
percent, the answer depended on which engine the underlying dictionary happened to
yield first. Twenty nine of those changed a digit. In a corpus of legal citations the
difference between \texttt{12-869} and \texttt{12369} is the difference between two
cases. The cause was mundane, and can affect any ensemble that uses first-encountered order as an
unstated tie break: \texttt{Counter.most\_common()} and \texttt{max()} resolve a tie by
returning whichever candidate was inserted first. We describe the four sites and their fix.
We also describe a fifth site outside the voter that the permutation test could not reach;
an external reviewer found it. The audited system treats refusal as an output rather than a
failure. We report its measured accuracy, and two reasons that number is an upper bound
rather than an estimate. Code, the reference standard and the replay fixture are public.
\end{abstract}

\section{Introduction}

Wrong characters are cheap. What costs money is a wrong character that nothing downstream
can tell apart from a right one.

This has become sharper with retrieval augmented generation. A pipeline that runs OCR
over a scanned filing and hands the result to a language model will reproduce a misread
docket number with the same fluency, and the same absence of hedging, as a correct one.
Unless the pipeline propagates confidence, disagreement or refusal, the model receives the
text as though every token were equally reliable. By the time it arrives, the distinction
between a corroborated reading and a guess has usually been discarded, and nothing
downstream can act on a difference it was never told about.

The system described here takes a narrow position on that problem. It does not try to
read scans more accurately than the best available engine. It tries to make every token in
its output carry a statement about how it was decided. Where its engines disagree it either
flags the reading or refuses to emit one, and lists the disagreement either way.

That design is not novel and we do not claim it is. Section~\ref{sec:related} sets out
how much of it is thirty years old. What we do claim is the finding in
Section~\ref{sec:audit}, and one methodological point: an ensemble whose tie breaks are
resolved by an unspecified rule is not reproducibly specified. The same inputs can yield
different outputs when insertion order changes, and none of the OCR ensemble papers we cite
reports a check for it.

\section{Related work}
\label{sec:related}

\paragraph{Voting across recognisers.} ROVER \cite{fiscus1997} aligns the output of
several speech recognisers into a word transition network and votes over it, optionally
weighting by each recogniser's confidence. The technique moved into document recognition
quickly. The ISRI toolkit at UNLV \cite{rice1994,nartker2005} shipped \texttt{synctext}
for alignment and \texttt{vote} for combination, and the documentation of \texttt{vote}
records that it applies majority voting with heuristics to break ties \cite{nartker2005}.
The problem in Section~\ref{sec:audit} was therefore visible in that documentation to
anyone who asked what those heuristics were, and we have found no report that anyone did.

Lund and Ringger \cite{lund2009} align several engines at the character level using an
admissible A$^{\star}$ heuristic and select from the resulting lattice with a dictionary,
reporting a 35.8 percent reduction in word error rate on a corpus of typewritten
communiqu\'es. Lund, Walker and Ringger \cite{lund2011} replace the selection step with a
discriminative model and report 24.6 percent relative improvement over the best single
engine. Handley \cite{handley1998} had already surveyed the space and separated
classifier level from string level combination.

\paragraph{Voting inside one engine.} Calamari \cite{wick2018} votes across models of a
single architecture trained on different folds, and reaches character error rates well
below anything we report, on cleaner material. Reul and colleagues \cite{reul2018} apply
the same idea to early printed books. This line of work is a genuine alternative to ours
rather than a predecessor: it gets its diversity from model instances, and we get ours
from differing architectures. Section~\ref{sec:independence} gives our reason for
preferring the second, which is a measurement rather than an argument.

\paragraph{Abstention.} Refusing to answer is not new either, and the literature is
considerably better developed than our use of it. Chow \cite{chow1957,chow1970}
formulated the reject option and derived the optimal threshold from the relative cost of
an error and a rejection. Modern selective prediction \cite{elyaniv2010,geifman2017}
reports a risk coverage curve rather than a single operating point. We now measure that
curve too (Section~\ref{sec:riskcov}); what we still lack is calibrated confidence, so
the curve is swept over the discrete acceptance rules rather than a continuous
threshold. Recent work scores OCR output by inter-model agreement entropy over several
vision language models \cite{zhang2025consensus}, and risk-controlled selective OCR has
been built around single vision language models \cite{geomrisk2026}. Both share our
premise that agreement or abstention carries information; neither votes across
independent classical engines, and the first selects a best output where we refuse. We
found no prior risk coverage measurement of a multi-engine OCR voting ensemble (nine
reviewing channels searched independently, 2026-08-09); that is a statement about what
the searches returned, not a claim of absence.

\paragraph{Ensembles as working software.} OCR-D's alignment step performs n-ary
alignment with majority or confidence voting as an operational component of a document
processing pipeline. Multi-engine OCR as usable software is not new, and a comparison
against that baseline is the most obvious experiment we have not run.

\section{The system}

Four voters read each page. One is the text layer already present in the PDF, whoever
produced it. Two are RapidOCR with a shared PP-OCRv6 detector and different recognisers.
The fourth is Tesseract~5.

Word boxes come from a frame engine chosen by a fixed preference order. Each voter's
reading is aligned to that frame, and a cascade of rules decides the token. The rules, in
order, are: exactly one engine produced anything; two or more produced the same
normalised string; the engines agree on the digits and differ on a separator; a numeric
token where a majority agree on the digit sequence; one candidate is in the lexicon and
the others are not; one candidate matches a citation pattern; no majority, in which case
the medoid is reported and flagged; and finally no defensible choice, in which case the
system refuses.

One property holds across all of them. Every rule returns a string that some engine
actually proposed, or nothing. No rule composes a new string from parts of several
readings. This is weaker than it sounds, because simple string level voters share it;
lattice and language model based reconciliation does not.

The output is a Markdown file whose header states, in fields that sum to the token count,
how many words were read identically by two or more independent engines, how many agreed
on digits but not punctuation, how many were settled by dictionary or pattern with no
second engine agreeing, and how many were not confirmed at all. The last group is written
to a separate file with page and coordinates.

\subsection{Independence is not the engine count}
\label{sec:independence}

Two of the four voters share a text detector. On one page set, 2{,}921 of 3{,}006 of
their boxes were identical. Counting their agreement as two independent confirmations
overstates the evidence, so the system counts independent opinions separately from
engines run, and flags any decision carried only by that correlated pair. On the sample
document, 4 of 856 confirmed words are in that state.

This is also the argument for architectural diversity over model instance diversity. The
question is not how many voters there are, it is how many ways of being wrong they
represent.

\section{The determinism audit}
\label{sec:audit}

\subsection{Method}

The decision function was extracted from the running system and a fixture was generated
by calling it, not by reimplementing it: for each of 29{,}853 token positions in a 41 page
United States immigration case index, the candidate readings and the resulting decision
were recorded. Deduplicating by candidate set gave 11{,}780 distinct decisions.

Each was then replayed under every permutation of the order in which candidates are
supplied, using \texttt{itertools.permutations}, and the set of returned triples
(chosen string, rule, backing engines) was required to have cardinality one.

\subsection{Result}

165 of 11{,}780 decisions, 1.40 percent, returned more than one distinct outcome
depending on ordering alone. Twenty nine changed a digit. Observed pairs include
\texttt{12-869} against \texttt{12369}, \texttt{1910} against \texttt{1940}, and
\texttt{11-335} against \texttt{11-336}.

Four sites were responsible. All four reduce to the same root cause. Python's
\texttt{Counter.\allowbreak most\_common()} and the built in \texttt{max()} are stable:
on a tie they return the candidate encountered first, and that order is the insertion
order of the mapping the engines were placed in. Nothing about that is a bug in Python. It is a
bug in using either function to make a decision that must be reproducible, without saying
what happens on a tie.

The fix in each case was to make the tie break a stated total order rather than an
accident: count first, then among the tied candidates select by a fixed key, and where no
key is defensible, fall through to a refusal instead of a choice. One consequence worth
reporting is that 33 of the 165 became refusals, that is, the corrected system declines
where the original silently picked.

\subsection{What the test could not see}

The permutation test covers the voting function. It does not cover the code that calls
it. A fifth ordering dependence was found later by an external reviewer in the selection
of the page frame, which was resolved alphabetically by engine name. The permutation test
could not reach it, because the fixture records decisions and not the frame the decisions
were made in.

A sixth was introduced by the author while fixing the fifth, in a metadata field that
records which engine supplied the frame, and was found by a further reviewer one commit
later. Both are worth reporting because they bound the claim: a permutation test over one
function proves a property of that function, and the natural reading of a green test as
``the system is deterministic'' is wrong.

More generally, permuting insertion order tests insertion order dependence, and not other
sources of nondeterminism such as set iteration, floating point summation order, or
engine version drift. Property based testing would cover more of that space than
exhaustive replay of recorded inputs does.

\subsection{Effect on the published accuracy}

The obvious objection is that the system's measured accuracy was produced by the code
this section shows to be nondeterministic. We re-scored the reference standard after the
four fixes. No labelled crop changed its verdict. The 165 affected decisions are
concentrated in the contested tail, and the reference standard, being stratified toward
contested cases, contains 60 crops of which none fell in that set. This is reassurance
rather than proof, and the correct statement is that the fixes did not move the measured
number on the sample available.

\section{Measurement}

\subsection{Against the pixels}

Sixty token positions were sampled from the index, stratified toward contested decisions,
and each was labelled by a human reading the page rendered at 900 dots per inch. Four
proved to have an unusable frame and are treated separately in
Section~\ref{sec:badbox}, leaving 56 scored positions.

\begin{table}[h]
\centering
\begin{tabular}{lrrc}
\toprule
voter & agrees with the page & counted & 95\% interval \\
\midrule
the PDF's own text layer & 46.4\% & 26 / 56 & [34.0, 59.3] \\
\texttt{rapidocr\_multi} & 78.6\% & 44 / 56 & [66.2, 87.3] \\
\texttt{rapidocr\_en} & 60.4\% & 32 / 53 & [46.9, 72.4] \\
\textbf{the ensemble} & \textbf{94.6\%} & \textbf{53 / 56} & \textbf{[85.4, 98.2]} \\
\bottomrule
\end{tabular}
\caption{Each voter and the ensemble against a human reading of the page. Intervals are
Wilson. Note that \texttt{rapidocr\_en} has a smaller denominator: it returned nothing at all
for three of the 56 positions, and a voter that abstains is not scored as wrong. The sample is
stratified toward contested decisions, so these are not document level accuracies and must not
be read as such.}
\label{tab:voters}
\end{table}

The comparison that matters in Table~\ref{tab:voters} is between the last row and the
best single row, and it is 16 points. The absolute values are not document accuracy: a
random sample of a document where the great majority of tokens are unanimous would spend
almost all of its labelling budget confirming easy words, which is why the sample is
stratified and why raw percentages from it are not extrapolated without weighting.

Weighted to the document with a Horvitz--Thompson estimator, accuracy is 98.65 percent
with three engines and 99.15 percent with four. An exact McNemar test on the discordant
pairs gives $p = 0.125$. We report the improvement because it was measured, and we report
the $p$ value because at this sample size the improvement cannot be distinguished from
chance.

\subsection{Four positions the measurement cannot see}
\label{sec:badbox}

Four of the sixty sampled positions have a word box, inherited from the PDF's text layer,
that does not correspond to a token. It cuts a glyph, or it frames typographic debris.
Every engine then reads inside the same wrong box, and every engine agrees.

This is the important negative result in the paper, and it is structural rather than
incidental. A reference standard built by rendering each token's own box and reading it
is wrong in exactly the same way the system is wrong, so it cannot detect the error.
Agreement based acceptance is blind to any failure the engines share, and a shared frame
is the most common way for them to share one. The headline accuracy figures are therefore
an upper bound on page level accuracy, not an estimate of it, and the gap between them is
unmeasured.

The same reasoning applies to any systematic error common to all engines: a font specific
confusion between \texttt{l} and \texttt{1}, ligature splitting, hyphenation at a line
break. Unanimity is evidence only to the extent that the voters fail independently, and
we have measured that independence for exactly one pair.

\subsection{The reference standard is an instrument}

The first version of the labels was read from a contact sheet whose vertical padding
caught the neighbouring line. Five of sixty labels were consequently taken from the wrong
line. They were found by a rule that is now part of the test suite and runs on every
invocation: if every engine agrees and the human disagrees, suspect the human. We
recommend it. An unchecked reference standard produces confident
errors, exactly as an uncalibrated instrument does.

\subsection{The curve, not one point}
\label{sec:riskcov}

Both measurements the first round of outside review ranked first are now in the
repository and recomputable by a stranger (\texttt{tests/test\_riskcov.py}): the risk
coverage curve over acceptance policies, and a single engine thresholded on its own
confidence to the same coverage. Coverage is a census and is exact; accuracy among
accepted tokens is a H\'ajek ratio over the 56 gold crops with post-stratification
weights by the rebuilt three-voice stratum. The draw-time inclusion probabilities no
longer exist -- the strata moved under the sample when two pipeline defects were fixed
-- so these are descriptive estimates under the shipped voter, not design-unbiased
inference from the original draw, and the same 56 crops served earlier rounds as a
development set, an adaptive reuse only a fresh sample on a second document can price.

\begin{table}[t]
\centering
\caption{Accepted-token accuracy against coverage: four of the fourteen measured
policies and the two single-engine baselines this stack can produce. Coverage is exact;
accuracy is a H\'ajek ratio over 56 stratified crops. The paired interval is the
bootstrap CI of the per-resample difference against the RapidOCR baseline, matched at
nearest coverage; marginal per-point intervals overlap everywhere.}
\label{tab:riskcov}
\footnotesize
\begin{tabular}{lrrl}
\toprule
policy & coverage & accuracy & honest uncertainty \\
\midrule
unanimity only & 77.3\,\% & 100\,\% & 0/7 errors; unweighted Wilson floor 64.6\,\% \\
VOTE$\geq$3, separators folded & 95.3\,\% & 99.4\,\% & CI [98.1, 100]; paired $\Delta$ [+0.6, +4.6] \\
full cascade & 97.6\,\% & 99.4\,\% & CI [98.2, 100]; paired $\Delta$ [+0.6, +4.6] \\
everything emitted & 98.9\,\% & 99.3\,\% & CI [97.9, 100]; Manski [97.7, 99.3] \\
\midrule
RapidOCR, line conf.\ $\tau{=}0.925$ & 98.4\,\% & 97.1\,\% & CI [94.7, 99.0] \\
Tesseract, word conf.\ $\tau{=}0.85$ & 80.6\,\% & 99.2\,\% & 20/21 crops; collapses below $\tau{=}0.8$ \\
\bottomrule
\end{tabular}
\end{table}

Three facts carry Table~\ref{tab:riskcov}. Measured accuracy does not collapse as
coverage rises: the later rules of the cascade add about fourteen points of coverage
while the point estimates move by less than 0.2 points, and the sample cannot
distinguish those policies from one another. The RapidOCR baseline is weak in an
instructive way: its confidence is line-level, every word inherits its line's score, so
thresholding barely sorts word errors (95.9\,\% accurate at 99\,\% coverage, 97.4\,\% at
14\,\%); against it the paired per-resample difference excludes zero at all fourteen
policies, in both matching directions. The fair challenger is Tesseract, whose
confidence is per word: at the coverages it can reach -- at most 83\,\% before its
accuracy collapses -- this sample cannot separate it from the ensemble. The measured
advantage of the ensemble over any single engine in this stack is therefore not accuracy
at low coverage; it is that no single engine's confidence delivers 95--99\,\% coverage
at comparable accuracy, and that is the coverage regime where a quotation tool operates.

The bounds, stated rather than implied. The four crops whose frame is unusable stay
excluded by the convention of Section~\ref{sec:limits}; counted as errors they cost the
full-coverage point half a percentage point and the unanimity point nothing. Removing
the lexicon entirely moves the headline accuracies by at most 0.01 and costs 0.3 points
of coverage, so the curve does not rest on it. Two strata hold a single crop and
contribute zero bootstrap variance; 560 accepted tokens (1.9\,\%) lie in strata with no
crop at all and enter only as Manski bands. The entire observed error set behind the
table is two crops, and both are printed in the repository's testing log.

\section{Limitations}
\label{sec:limits}

We list these in the order we would attack them.

One document, one genre, one language, one scanning lineage from 2002. Nothing here shows
that any number transfers.

The risk coverage curve of Section~\ref{sec:riskcov} is measured over discrete
policies, not a continuous threshold, and its accuracy axis rests on 56 crops with two
strata behind a single crop each. The single-engine comparison is honest but asymmetric
in both directions at once: RapidOCR's line-level confidence understates what a
well-calibrated single engine could do, and Tesseract's word-level confidence cannot
reach the coverage where the comparison matters. A single engine with calibrated
word-level confidence at high coverage does not exist in this stack; testing against one
elsewhere is an open task, and if such an engine dominated the curve, the premise of the
system would be weaker than we have argued.

No decision level reference standard. The 56 labelled positions measure whether the
chosen string matches the page. They do not measure, per rule, how often a rule that
fires is right, except at sample sizes where the intervals span most of the unit
interval.

No downstream experiment. The motivating claim is that enumerated uncertainty reduces
fabrication when the text reaches a language model. That claim is untested here.

No comparison against an external ensemble. OCR-D's alignment and voting component is
available and would be the natural baseline.

The lexicon used by one rule was not audited for overlap with the test document. If it
was derived from material that includes it, that rule's accuracy is overstated.

\section{Tools used in preparing this work}

The software described here was written with substantial assistance from a large language
model, and so was this report. The measurements, the fixture, the reference standard and the
audit in Section~\ref{sec:audit} were produced by executing code and recording what it
returned, not by asking a model what it thought the answer was. Several of the defects
described in Sections~\ref{sec:audit} and \ref{sec:limits}, including the fifth ordering
dependence and an error in an earlier version of the accuracy accounting, were found by
independent language models asked to review the artefact adversarially.

The author is responsible for the entire contents, including every citation, and has checked
them.

\section{Availability}

Code, the five test documents with their SHA-256 checksums and government source URLs,
the replay fixture of 11{,}780 decisions, and the 60 position reference standard with its
labels are at \url{https://github.com/igorsaevets/scanquorum} under the MIT licence. The
accuracy figures in Table~\ref{tab:voters} are recomputed by a test in that repository
rather than quoted from a notebook.

\begin{thebibliography}{99}

\bibitem{chow1957} C.~K. Chow. An optimum character recognition system using decision
functions. \emph{IRE Transactions on Electronic Computers}, EC-6(4):247--254, 1957.

\bibitem{chow1970} C.~K. Chow. On optimum recognition error and reject tradeoff.
\emph{IEEE Transactions on Information Theory}, 16(1):41--46, 1970.

\bibitem{elyaniv2010} R.~El-Yaniv and Y.~Wiener. On the foundations of noise-free
selective classification. \emph{Journal of Machine Learning Research}, 11:1605--1641,
2010.

\bibitem{fiscus1997} J.~G. Fiscus. A post-processing system to yield reduced word error
rates: Recognizer Output Voting Error Reduction (ROVER). In \emph{IEEE Workshop on
Automatic Speech Recognition and Understanding}, pages 347--354, 1997.

\bibitem{geifman2017} Y.~Geifman and R.~El-Yaniv. Selective classification for deep
neural networks. In \emph{Advances in Neural Information Processing Systems 30}, pages
4878--4887, 2017.

\bibitem{geomrisk2026} W.~Gong, Z.~Lu, M.~Chen, Y.~Zuo, X.~He, and W.~Fan. Geometric
risk control for vision-language model OCR. arXiv:2603.19790, 2026.

\bibitem{handley1998} J.~C. Handley. Improving OCR accuracy through combination: a
survey. In \emph{IEEE International Conference on Systems, Man, and Cybernetics}, volume
5, pages 4330--4333, 1998.

\bibitem{lund2009} W.~B. Lund and E.~K. Ringger. Improving optical character recognition
through efficient multiple system alignment. In \emph{Joint Conference on Digital
Libraries}, pages 231--240, 2009.

\bibitem{lund2011} W.~B. Lund, D.~D. Walker, and E.~K. Ringger. Progressive alignment and
discriminative error correction for multiple OCR engines. In \emph{International
Conference on Document Analysis and Recognition}, pages 764--768, 2011.

\bibitem{nartker2005} T.~A. Nartker, S.~V. Rice, and S.~E. Lumos. Software tools and test
data for research and testing of page-reading OCR systems. In \emph{Document Recognition
and Retrieval XII}, volume 5676 of \emph{Proceedings of SPIE}, pages 37--47, 2005.

\bibitem{reul2018} C.~Reul, U.~Springmann, C.~Wick, and F.~Puppe. Improving OCR accuracy
on early printed books by utilizing cross fold training and voting. In \emph{13th IAPR
International Workshop on Document Analysis Systems}, pages 423--428, 2018.

\bibitem{rice1994} S.~V. Rice, J.~Kanai, and T.~A. Nartker. An algorithm for matching
OCR-generated text strings. \emph{International Journal of Pattern Recognition and
Artificial Intelligence}, 8(5):1259--1268, 1994.

\bibitem{wick2018} C.~Wick, C.~Reul, and F.~Puppe. Calamari: a high-performance
Tensorflow-based deep learning package for optical character recognition.
arXiv:1807.02004, 2018.

\bibitem{zhang2025consensus} Y.~Zhang, T.~Liang, X.~Huang, E.~Cui, G.~Wang, X.~Guo,
C.~Li, and G.~Liu. Consensus entropy: harnessing multi-VLM agreement for self-verifying
and self-improving OCR. arXiv:2504.11101, 2025.

\end{thebibliography}

\end{document}
